\chapter{Introduction}
There are one set of people who wants to pursue higher education degrees due to their passion and interest. For these people, it makes sense to get enrol in the relevant university courses and pursue their passion. On the other hand, there is another set of people, who wants to obtain these degrees only to get a specific job title, a specific job role, or a position that get them more money. For this set, university degrees are not the only option, there are many alternatives which can also result in the same outcomes.

The most obvious example is in the field of Data Science and Analytics. In recent years, university degrees such as "Masters in Data Science" or "Masters in Analytics" are sought as one of the must-haves to enter into this field. It is not astonishing that Data Scientist is one of the fastest-growing job titles across the globe and the demand for skilled data scientists is increasing. This has given the universities an option to attract students and make immense money. Several universities have started dedicated degree courses specialized in data science and analytics. Those want to become a data scientist or to switch from another profession to data science profession are now strongly considering these university degrees as the only pathway.

But these university courses are not easy to get in and affordable for everyone. These degrees don’t come for free, tuition fees can be exorbitant and can range anything from USD 30,000 to USD 100,000. And that doesn’t include the actual cost of living. Many consider applying for student loans but they add a huge lump sum to the existing mountain of debts. A common myth is that the earning potential for those with postgraduate qualifications is higher but of course, there is no guarantee that one will get a stable job at the end of it. Additionally, Pursuing a university’s higher degree takes anything from one to three years, depending on different factors. This can seem like a long time, especially when the fellow peers are getting started on their careers, while one is still studying.

The question of interest here is - \textbf{does one need to get that expensive higher education degree, do they create a difference from those who do not have university degrees?}. Some resources online also suggest that one can get the depth of knowledge, variety of skills and learn something new. But again, \textbf{is it possible to get the same skills, same profile, or even better compensation without such degrees?}.

\section{Source of the Data}

\href{www.kaggle.com}{Kaggle} conducted their Annual Data Science survey and it was full of interesting questions. Participants of this survey were asked different questions about their demographics, profiles, companies, what they use etc. I analysed this data intending to dig deeper into the profiles of people who completed the university degrees to learn data science and those who did not. The focus of the study is to identify if the working data scientists with official higher education degrees differ significantly from the other group. The analysis and storyline are segmented according to different factors.

\note{The respondents who were "students" and "not employed" were removed for analysis. The two groups were selected based on the respondent's choice if they completed the university degrees to learn data science or not.}
